%% kapa: exam=1 round=25 year=2014 period=1 form=A subject=economics subject_ko=경제원론 seq=18 number=58
%% source: Q-Net 감정평가사 기출문제 (cst00309) artlSeq=5101207 file=extracted/1008/5101207/2014년 제25회 감정평가사 1차시험 기출문제/1교시(A형).hwp
%% conversion: kordoc 4.24.51 -> tools/gen1.py (auto); review=proofread-sonnet; flags=none
\begin{question}{58}
$U= \sqrt { Y }$의 효용함수를 갖는 소비자가 100만원의 가치가 있는 자전거를 소유하고 있다. 자전거의 도난확률이 0.5일 때 다음 중 옳지 않은 것은? (단, $Y$는 재화가치이다.)
\begin{choices}
  \item 위험한 기회를 다른 사람에게 전가할 때 지급할 최대 추가보상액은 50만원이다.
  \item 현재 이 소비자의 기대효용수준은 500이다.
  \item 손실액 전액을 보상해주는 보험의 경우 공정한 보험료는 50만원이다.
  \item 손실액 전액을 보상해주는 보험에 대해 이 소비자는 최대 75만원까지 지불할 용의가 있다.
  \item 이 소비자는 위험기피자이다.
\end{choices}
\end{question}
\answer{1}
